\chapter{Direction des énergies et Centre de Grenoble}


\section{Origines et contributions à la recherche sur les énergies}

    \begin{figure}[H]
        \centering
        % Le contenu du diagramme est défini dans le fichier suivant
        \fbox{\begin{tikzpicture}[
    main/.style={text centered, text width=10cm, fill=red!20, draw, rectangle, rounded corners, font=\bfseries\Large},
    direction/.style={text centered, text width=3cm, fill=gray!50, draw, rectangle, rounded corners},
    institut/.style={grow=down, xshift=4cm, text centered, text width=5cm, draw, rectangle, rounded corners, edge from parent path={(\tikzparentnode.225) |- (\tikzchildnode.west)}},
    first/.style    ={level distance=15ex},
    second/.style   ={level distance=30ex},
    third/.style    ={level distance=45ex},
    fourth/.style   ={level distance=60ex},
    fifth/.style    ={level distance=75ex},
    level 1/.style={sibling distance=12.5em, level distance=8em}]

    \node[main, anchor=south]{Direction Générale du CEA}
    [edge from parent fork down]
    %
    child{node (d1) [direction] {\textbf Direction \\ des \textbf Énergies}
    child[institut,first] {node {\textbf{I}nstitut de {R}echerche\\ et d'{É}{\textbf t}ude{\textbf s} en {É}conomie\\ de l'\textbf{É}nergie}}
    child[institut,second] {node {\textbf{L}aboratoire d'\textbf{I}nnovation\\ pour les \textbf{T}echnologies des\\ \textbf{É}nergies {N}ouvelles et\\ les \textbf{N}anomatériaux}}
    child[institut,third] {node {\textbf{I}nstitut des \textbf{S}ciences\\ \textbf{et} {T}echnologies pour une\\ \textbf{É}conomie \textbf{C}irculaire\\ des {é}nergies bas carbone}}
    child[institut,fourth] {node {\textbf{I}nstitut de \textbf{Re}cherche\\ pour les \textbf{S}ystèmes \textbf{N}ucléaires\\ pour la production d'\textbf{É}nergie bas carbone}}
    child[institut,fifth] {node {\textbf{I}nstitut des \textbf{S}ciences\\ \textbf{A}ppliquées et de la \textbf{S}imulation\\ pour les {É}nergies carbone}}
    }
    %
    child{node (d2) [direction] {\textbf Direction \\ des \textbf{A}pplications \textbf{M}ilitaires}}
    %
    child{node (d3) [direction] {\textbf Direction \\ de la \textbf{R}echerche \textbf{T}echnologique}}
    %
    child{node (d4) [direction] {\textbf Direction \\ de la \textbf{ R}echerche \textbf{F}ondamentale}};

\end{tikzpicture}}
        \caption{Organigramme simplifié des Directions Opérationnelles du CEA}
        \label{fig:organigramme-CEA}
    \end{figure}

    \begin{figure}[H]
        \centering
        % Le contenu du diagramme est défini dans le fichier suivant
        \fbox{\begin{tikzpicture}[
    main/.style={text centered, text width=10cm, fill=red!50, draw, rectangle, rounded corners, font=\bfseries\Large},
    direction/.style={text centered, text width=3cm, fill=orange!20, draw, rectangle, rounded corners, font=\large},
    institut/.style={grow=down, xshift=3cm, text centered, text width=4cm, draw, rectangle, rounded corners, edge from parent path={(\tikzparentnode.225) |- (\tikzchildnode.west)}},
    site/.style={level distance = 6ex, xshift=4.5cm, text centered, text width=5cm, draw, rectangle, rounded corners, edge from parent path={(\tikzparentnode.225) |- (\tikzchildnode.west)}},
    first/.style    ={level distance=12ex},
    second/.style   ={level distance=24ex},
    third/.style    ={level distance=36ex},
    fourth/.style   ={level distance=48ex},
    fifth/.style    ={level distance=60ex},
    level 1/.style={sibling distance=30em, level distance=5em}]

    \node[main, anchor=south]{Direction Générale du CEA}
    [edge from parent fork down]
    %
    child{node (d1) [direction] {Centres civils}
        child[institut,first]  {node {Centre de Cadarache}}
        child[institut,second] {node {Centre de Grenoble}
            child[site]  {node {\textbf Institut \textbf National de l'\textbf Énergie \textbf Solaire (Chambéry)}}
        }
        child[institut,third]  {node {Centre de Marcoule}}
        child[institut,fourth] {node {Centre de Paris-Saclay}}
    }
    %
    child{node (d2) [direction] {Centres Militaires}
        child[institut,first]  {node {Centre du Cesta}}
        child[institut,second] {node {Centre de Gramat}}
        child[institut,third]  {node {Centre du Ripault}}
        child[institut,fourth] {node {Centre de Valduc}}
        child[institut,fifth] {node {Centre DAM Île-de-France}}
    };
\end{tikzpicture}}
        \caption{Organigramme simplifié des Centres du CEA}
        \label{fig:organigramme-Centres}
    \end{figure}

    - Prix Nobels : ... + Anne L'Huillier en 2023

    Citons quelques chiffres-clés du CEA\cite{CEA_prez_generale} pour conclure cette introduction :
    \begin{itemize}
        \item
    \end{itemize}

    \subsection{Direction des énergies}

    \cite{CEA_prez_DES}

% Focus sur Grenoble + histoire
\section{Le Centre CEA de Grenoble : entre diversité des domaines de recherche et résilience historique}

Conçu à l'origine pour… Le Centre de Grenoble a su se réinventer,

Incendie des salles blanches : problème pour l'image de marque et la crédibilité mais surpassement de ces difficultés

\ldots\footnote{Il convient de préciser que l'alternance rapportée ne se situe pas réellement à Grenoble mais sur le site de l'INES en Savoie, rattaché administrativement au Centre de Grenoble. Ainsi, l'INES fait partie des nombreuses autres antennes réparties sur tout le territoire français permettant de faire du CEA un acteur national de premier plan en matière de recherche technologique.}

\cite{CEA_Gre}


